\documentclass[11pt,a4paper]{article}
\usepackage[margin=1in]{geometry}
\usepackage{hyperref}
\usepackage{enumitem}
\usepackage{graphicx}
\usepackage{xcolor}
\usepackage{booktabs}

% -----------------------------------------------------
% bvraghav-tiet-ucs503p-article.sty
% -----------------------------------------------------

\makeatletter \newcommand{\BvrSetLabInstructor}[1]{%
  \gdef\Bvr@LabInstructor{#1}}
\newcommand{\Bvr@LabInstructor}{Dr. Raghav %
  B. Venkataramaiyer}
\newcommand{\BvrLabInstructorValue}{\Bvr@LabInstructor}
\makeatother

\usepackage{titling}
\pretitle{\begin{center}\bfseries\fontsize{18bp}{18bp}\selectfont}
\makeatletter
\newcommand{\BvrSubtitle}[1]{%
  \posttitle{%
    \par\end{center}
    \begin{center}\large#1\end{center}
    \vskip 0.5em}%
}
\makeatother

\newcommand{\BvrHint}[1]{%
  {\normalfont\normalsize\itshape\hfill #1}}

% -----------------------------------------------------

\title{SportSathi: An AI-Based Fitness Screening and Sports Talent Guidance Platform for Schools}

\BvrSetLabInstructor{Dr. Raghav B. Venkataramaiyer}

\BvrSubtitle{%
  Project Report --- Prototype Stage, for UCS503P  \\
  Submitted to: \BvrLabInstructorValue \\ \bfseries
  Thapar Institute of Engineering and Technology%
}

\author{
  Kiratpreet Kaur, CSED%
  \thanks{Roll No: \texttt{1024030402}, Batch: \texttt{3C31}, %
    Email: \texttt{kkaur1\_be24\@thapar.edu}}
  \and
  Prisha Goyal, CSED%
  \thanks{Roll No: \texttt{1024030404}, Batch: \texttt{3C31}, %
    Email: \texttt{pgoyal1\_be24\@thapar.edu}}
}

\date{\today}

\begin{document}
\maketitle

\section{Purpose of this stage}
This report captures the state of SportSathi midway
between the initial proposal and the final report. Its
purpose is narrower than the final report: to show that the
riskiest technical assumption in the proposal --- that an
explainable classifier can be trained honestly on available
student fitness data --- has been tested, and to record what
was learned before the remaining workflow (coach approval,
dashboards, reporting) was built out.

\section{What was prototyped}
\begin{itemize}[leftmargin=0.7cm]
    \item A public student physical-education performance dataset (500 records, 20 attributes) was identified and adopted as the Tier~1 training/evaluation dataset, in place of collecting real school data, consistent with the proposal's privacy constraints.
    \item A first baseline classifier (Decision Tree) was trained and evaluated in isolation, ahead of any backend integration, exactly as the proposal's time-to-value heuristic specified.
    \item The assessment-entry form and BMI auto-calculation were scaffolded on the frontend.
\end{itemize}

\section{Key finding: a leakage problem the proposal did not anticipate}
The first baseline run included a derived
\texttt{Overall\_PE\_Performance\_Score} column as a model
input and produced a suspicious 98\% accuracy. Feature
importance analysis showed the model relying almost
entirely on that single derived column and on attendance,
while every genuine physical fitness score (strength,
endurance, flexibility, speed-agility) had zero importance
--- a clear sign the derived column already encoded the
answer.

This was corrected by excluding all derived/summary columns
(\texttt{Overall\_PE\_Performance\_Score},
\texttt{Improvement\_Rate}, \texttt{Final\_Grade}) from the
feature set. Retraining on only genuine, independently
collected attributes produced a more credible 81--86\%
accuracy (Decision Tree: 86\%, Random Forest: 81\%).

\textbf{This finding is the main output of the prototype
stage.} It confirms the proposal's ``validate in isolation
before integration'' heuristic was the right call: the
problem was caught while the model was still being tested
on its own, before it was wired into any user-facing
workflow.

\section{Secondary finding: attendance as a strong predictor}
Class-wise average attendance showed a clean, monotonic
separation across performance categories (High Performer:
94.2\%, Average: 89.8\%, Low: 84.7\%), stronger than any
single raw physical-fitness score at separating the
classes. This was not anticipated in the proposal and is
carried forward as a candidate insight to validate further
in the final stage.

\section{Known limitation carried forward}
The dataset's target classes are heavily imbalanced (245
Low Performer, 243 Average Performer, 12 High Performer).
Both the Decision Tree and Random Forest baselines failed
to correctly predict any High Performer example. This is
flagged here, before further work, as a limitation to be
explicitly addressed (or explicitly documented as
unresolved) in the final report, rather than discovered
late.

\section{Status against the proposal's initial deliverable list}
\begin{center}
\begin{tabular}{@{}p{6.5cm}p{4.5cm}@{}}
\toprule
Proposed initial deliverable & Status at prototype stage \\
\midrule
Authentication and role-based access & Not yet started \\
Tier 1 assessment entry + BMI calc. & Frontend scaffolded \\
Trained, evaluated classifier & Baseline complete; leakage found and corrected \\
Coach approval/override workflow & Not yet started \\
CI: build, tests, linting & Not yet started \\
CD to staging & Not yet started \\
\bottomrule
\end{tabular}
\end{center}

\section{Revised plan for the final stage}
\begin{itemize}[leftmargin=0.7cm]
    \item Integrate the corrected classifier into the backend, behind a coach-approval gate, before adding any further dashboard features.
    \item Apply class-imbalance handling (SMOTE, class weighting) to the classifier and report per-class metrics, not just overall accuracy, in the final evaluation.
    \item Prioritise the coach-approval workflow and student/parent visibility gating over CI/CD if a time trade-off becomes necessary, since the approval gate is the project's core safety property.
    \item Treat sport-suitability recommendation logic and the full relational database migration as stretch goals for the final stage, to be reported honestly as complete or incomplete rather than assumed.
\end{itemize}

\section{Summary}
The prototype stage validated the proposal's central
technical risk --- that an explainable, honestly-evaluated
fitness classifier is achievable on available data --- and,
in doing so, surfaced a data-leakage problem the proposal
had not anticipated. Catching this early, during isolated
model testing rather than after integration, is treated
here as evidence that the proposal's time-to-value approach
is working as intended. The remaining workflow components
(coach approval, dashboards, CI/CD) are carried into the
final stage as the primary remaining scope.

\end{document}
